\subsection{Derivations in the algebra of formal power series}

Let I be a set, E an associative, commutative and unital linearly topologized, separated and complete A-algebra, and \(\mathbf{x} = (x_{i})_{i \in I}\) a family of elements of E satisfying conditions a) and b) of Prop. 4 (IV, p.~28). Let \(\varphi\) be the continuous homomorphism \(u \mapsto u(x)\) of A{[}{[}I{]}{]} into E; it equips E with an A{[}{[}I{]}{]}-module structure. By III, p.~552, an A-derivation D of A{[}{[}I{]}{]} into the A{[}{[}I{]}{]}-module E is thus an A-linear mapping D: A{[}{[}I{]}{]} → E satisfying the relation

\[
\mathrm{D} (u v) = u (\mathbf {x}) \cdot \mathrm{D} (v) + \mathrm{D} (u) \cdot v (\mathbf {x})\tag{15}
\]

for u, v in A{[}{[}I{]}{]}.

PROPOSITION 7. --- Let \((y_{i})_{i\in I}\) be a family of elements of E. When I is infinite, we assume that \(y_{i}\) tends to 0 along the filter of complements of finite subsets of I. There exists then a unique continuous A-derivation D of A{[}{[}I{]}{]} into the A{[}{[}I{]}{]}-module E such that \(\mathrm{D}(\mathbf{X}_{i})=y_{i}\) for all \(i\in I\) . We have

\[
\mathrm{D} (u) = \sum_ {i \in I} (\mathrm{D} _ {i} u) (\mathbf {x}) \cdot y _ {i} \quad (u \in \mathrm{A} [ [ I ] ])  .\tag{16}
\]

Since 0 admits in E a fundamental system of neighbourhoods consisting of ideals, the family \(((\mathrm{D}_i u)(\mathbf{x}) \cdot y_i)_{i \in I}\) is summable in E for all \(u \in A[[I]]\) (Gen.~Top., III, p.~263, Cor. 2). Formula (16) thus defines an A-linear mapping \(\mathbf{D}: A[[I]] \to E\) . We leave the reader to verify that D is a continuous derivation. Let \(\mathbf{D}_1\) be a continuous A-derivation of A{[}{[}I{]}{]} into E, such that \(\mathbf{D}_1(\mathbf{X}_i) = y_i\) for all \(i \in I\) . The kernel of the continuous derivation \(\mathbf{D} - \mathbf{D}_1\) is a closed subalgebra B of A{[}{[}I{]}{]} containing 1 and the indeterminates \(X_{i}\) . Since the polynomial algebra \(\mathbf{A}[(\mathbf{X}_{i})_{i\in I}]\) is dense in A{[}{[}I{]}{]}, we have \(B = A[[I]]\) and so \(D_{1} = D\) .

COROLLARY 1. --- Let \(\Delta\) be a continuous derivation of the A-algebra E. For every formal power series \(u \in \mathbf{A}[[\mathrm{I}]]\) the family \(((\mathrm{D}_i u)(\mathbf{x}) \cdot \Delta x_i)_{i \in I}\) is summable and we have

\[
\Delta (u (\mathbf {x})) = \sum_ {i \in I} \left(\mathrm{D} _ {i} u\right) (\mathbf {x}) \cdot \Delta x _ {i}.\tag{17}
\]

This follows from Prop. 7 because the mapping \(u \mapsto \Delta(u(\mathbf{x}))\) is a continuous derivation of \(\mathbf{A}[[\mathbf{I}]]\) into the \(\mathbf{A}[[\mathbf{I}]]\) -module \(\mathbf{E}\) .

COROLLARY 2. --- The derivation \(D_{i}\) is the unique continuous derivation of the A-algebra A{[}{[}I{]}{]} such that

\[
\mathrm{D} _ {i} (\mathrm{X} _ {i}) = 1, \quad \mathrm{D} _ {i} (\mathrm{X} _ {j}) = 0 \quad f o r \quad j \neq i.\tag{18}
\]

This follows from Cor. 1.

COROLLARY 3. --- Let \(f \in \mathbf{A}[[\mathbf{X}_1, \ldots, \mathbf{X}_p]]\) and \(g_i \in \mathbf{A}[[\mathbf{Y}_1, \ldots, \mathbf{Y}_q]]\) for \(1 \leqslant i \leqslant p\) . Suppose that for \(1 \leqslant i \leqslant p\) the constant term of \(g_i\) is zero, and put \(h = f(g_1, \ldots, g_p)\) . Then for \(1 \leqslant j \leqslant q\) we have

\[
\frac {\partial h}{\partial \mathbf {Y} _ {j}} = \sum_ {i = 1} ^ {p} \mathrm{D} _ {i} f (g _ {1}, \dots , g _ {p}) \cdot \frac {\partial g _ {i}}{\partial \mathbf {Y} _ {j}}.\tag{19}
\]

This is the special case \(\mathbf{E} = \mathbf{A}[[\mathbf{Y}_1,\dots ,\mathbf{Y}_q]]\) , \(x_{i} = q_{i}\) and \(\Delta = \partial /\partial \mathbf{Y}_j\) of Cor. 1.

PROPOSITION 8. --- Let \(\mathbf{X} = (\mathbf{X}_i)_{i\in I}\) be a finite family of indeterminates.

\begin{enumerate}
\def\labelenumi{(\roman{enumi})}
\item
  Every derivation of the ring of formal power series A{[}{[}X{]}{]} is continuous.
\item
  Every derivation of the polynomial ring A{[}X{]} into the ring of formal power series A{[}{[}X{]}{]} extends in a unique fashion to a derivation of the ring A{[}{[}X{]}{]}.
\item
  The family \((\mathbf{D}_i)_{i\in I}\) is a basis of the \(\mathbf{A}[[\mathbf{X}]]\) -module of \(\mathbf{A}\) -derivations of \(\mathbf{A}[[\mathbf{X}]]\) into itself.
\end{enumerate}

Let \(\mathfrak{b}_n\) be the set of all formal power series of order \(\geqslant n\) . It is clear that \(\mathfrak{b}_n\) is an ideal in the ring \(A[[X]]\) , generated by the monomials of degree \(n\) . Hence \(\mathfrak{b}_n\) consists of finite sums of products of \(n\) formal power series without constant terms; if \(D\) is a derivation of \(A[[X]]\) , we have

\[
\mathrm{D} (f _ {1} \dots f _ {n}) = \sum_ {i = 1} ^ {n} f _ {1} \dots f _ {i - 1} \mathrm{D} (f _ {i}) f _ {i + 1} \dots f _ {n},
\]

whence it follows at once that \(\mathbf{Db}_{n} \subset \mathbf{b}_{n-1}\) for \(n \geqslant 1\) . Since the sequence \((\mathfrak{b}_n)_{n \geqslant 0}\) is a fundamental system of neighbourhoods of 0 in A{[}{[}X{]}{]} (IV, p.~26 remark), D is continuous and (i) is proved.

Let \(\Delta\) be a derivation of A{[}X{]} into A{[}{[}X{]}{]}. Arguing as before, we can show that \(\Delta(h)\) belongs to \(\mathfrak{b}_{n-1}\) for every homogeneous polynomial \(h\) of degree \(n \geqslant 1\) . Now let \(u \in \mathbf{A}[[\mathbf{X}]]\) and let \(u_n\) be the homogeneous component of degree \(n\) of \(u\) . Since \(\Delta(u_n) \in \mathfrak{b}_{n-1}\) for \(n \geqslant 1\) , the family \((\Delta(u_n))_{n \geqslant 0}\) is summable in A{[}{[}X{]}{]} and we can define a derivation D of A{[}{[}X{]}{]} into itself by

\[
\mathrm{D} (u) = \sum_ {n \geqslant 0} \Delta \left(u _ {n}\right).
\]

We have \(D(\mathfrak{b}_n) \subset \mathfrak{b}_{n-1}\) , hence \(D\) is a continuous endomorphism of the additive group of \(A[[X]]\) . The mapping \(\Phi: (u, v) \mapsto D(uv) - uD(v) - D(u)v\) of \(A[[X]] \times A[[X]]\) into \(A[[X]]\) is continuous and zero on \(A[X] \times A[X]\) . Since \(A[X]\) is dense in \(A[[X]]\) , we have \(\Phi = 0\) ; in other words, \(D\) is a derivation of \(A[[X]]\) into itself, extending \(\Delta\) .

Finally, A{[}X{]} is dense in A{[}{[}X{]}{]} and every derivation of A{[}{[}X{]}{]} is continuous by (i); hence there exists a unique extension of \(\Delta\) to a derivation of A{[}{[}X{]}{]}. This proves (ii).

It remains to prove (iii). Formula (18) (IV, p.~33) shows that the family \((\mathbf{D}_i)_{i\in I}\) is linearly independent over \(\mathbf{A}[[\mathbf{X}]]\) , and formula (16) (IV, p.~32), applied in the case \(\mathbf{E} = \mathbf{A}[[\mathbf{X}]]\) , shows that every A-derivation is a linear combination of the \(\mathbf{D}_i\) with coefficients in \(\mathbf{A}[[\mathbf{X}]]\) .

PROPOSITION 9. --- Let \((u_{\lambda})_{\lambda \in L}\) be a summable family of elements of A{[}{[}I{]}{]} without constant term and D a continuous derivation of the A-algebra A{[}{[}I{]}{]}. If \(f = \prod_{\lambda \in L} (1 + u_{\lambda})\) (IV, p.~27, Prop. 2), then the family \((D u_{\lambda} / (1 + u_{\lambda}))_{\lambda \in L}\) is

summable and we have

\[
\mathrm{D} (f) / f = \sum_ {\lambda \in L} \mathrm{D} (u _ {\lambda}) / (1 + u _ {\lambda}).\tag{20}
\]

If \(g\) and \(h\) are two invertible elements of \(\mathbf{A}[[\mathbf{I}]]\) , then

\[
\mathrm{D} (g h) = h. \mathrm{D} g + g. \mathrm{D} h
\]

whence on division by gh,

\[
\mathrm{D} (g h) / g h = \mathrm{D} (g) / g + \mathrm{D} (h) / h.\tag{21}
\]

For every finite subset M of L put \(f_{M} = \prod_{\lambda \in M} (1 + u_{\lambda})\) . From (21) we deduce by induction on Card M the relation

\[
\mathrm{D} \left(f _ {\mathrm{M}}\right) / f _ {\mathrm{M}} = \sum_ {\lambda \in \mathrm{M}} \mathrm{D} \left(u _ {\lambda}\right) / \left(1 + u _ {\lambda}\right).\tag{22}
\]

This proves Prop. 9 when L is finite. Now suppose that L is infinite and write F for the filtered ordered set of finite subsets of L. We have \(\lim_{\mathfrak{F}} f_{M} = f\) , and hence (IV, p.~30 remark)

\[
\mathrm{D} (f) / f = \lim _ {\mathfrak {F}} \mathrm{D} (f _ {\mathrm{M}}) / f _ {\mathrm{M}}.
\]

Now Prop. 9 follows by passage to the limit in (22).

